\documentclass[11pt]{article}
\usepackage[margin=1.2in]{geometry}
\usepackage{amsmath,amssymb,amsthm,mathtools}
\usepackage{hyperref}

\newtheorem{theorem}{Theorem}
\newtheorem{proposition}[theorem]{Proposition}
\newtheorem{lemma}[theorem]{Lemma}
\theoremstyle{remark}
\newtheorem{remark}[theorem]{Remark}

\newcommand{\PP}{\mathbb{P}}
\newcommand{\Dcal}{\mathrm{D}}
\newcommand{\coh}{\mathrm{coh}}
\newcommand{\grmod}{\mathrm{grmod}}
\newcommand{\stgrmod}{\underline{\mathrm{grmod}}}
\newcommand{\perf}{\mathrm{perf}}
\newcommand{\Hom}{\operatorname{Hom}}
\newcommand{\Sym}{\operatorname{Sym}}
\renewcommand{\k}{\mathbf{k}}

\title{The BGG correspondence: strategy of the classical proof}
\author{}
\date{}

\begin{document}
\maketitle

\section{Statement}

Let $V$ be an $(n+1)$-dimensional vector space over a field $\k$, $W=V^*$, and
\[
S=\Sym(V),\qquad E=\Lambda(W),
\]
so that $\PP^n=\PP(W)$ has homogeneous coordinate ring $S$. Grade $S$ by $\deg V=1$ and $E$ by $\deg W=-1$. The algebra $E$ is a graded Frobenius algebra with $\omega_E\cong\Hom_\k(E,\k)\cong E(-n-1)$ up to shift convention, so finitely generated free $E$-modules are both projective and injective.

\begin{theorem}[Bernstein--Gelfand--Gelfand, 1978]
There is an equivalence of triangulated categories
\[
\Dcal^b(\coh \PP^n)\;\simeq\;\stgrmod E\;\simeq\;\Dcal^b(\grmod E)/\perf(E),
\]
where $\stgrmod E$ is the stable category of finitely generated graded $E$-modules.
\end{theorem}

The proof has three ingredients: Koszul duality between $S$ and $E$, a comparison of the two Serre/Verdier quotients, and (for explicit inverses) the Tate resolution and Beilinson's resolution of the diagonal.

\section{Step 1: Koszul duality}

For a graded $S$-module $M$ and a graded $E$-module $P$ define the linear complexes
\begin{align*}
\mathbf{R}(M)&:\ \cdots\to \Hom_\k(E,M_d)\xrightarrow{\ \sum e_i\otimes x_i\ }\Hom_\k(E,M_{d+1})\to\cdots &&\text{(free $E$-modules)},\\
\mathbf{L}(P)&:\ \cdots\to S\otimes_\k P_d\xrightarrow{\ \sum x_i\otimes e_i\ }S\otimes_\k P_{d+1}\to\cdots &&\text{(free $S$-modules)},
\end{align*}
where $x_i$, $e_i$ are dual bases of $V$, $W$. The module structure on one side becomes the differential on the other.

\begin{theorem}[BGG, Koszul duality]
$\mathbf{R}$ and $\mathbf{L}$ induce mutually inverse equivalences
\[
\Dcal^b(\grmod S)\;\simeq\;\Dcal^b(\grmod E).
\]
\end{theorem}

\emph{Idea.} The Koszul complex $S\otimes\Lambda(V)$ resolves $\k$ over $S$, and dually over $E$. Hence $\mathbf{R}(S(d))$ is, up to shift, the residue field $\k$ over $E$, while $\mathbf{R}(\k)$ is a free module ($\omega_E$ up to shift). Since both categories are generated by these objects, and $\Hom$ spaces between generators match (Koszul property), $\mathbf{R}$ is fully faithful and essentially surjective. (Shifts and twists depend on conventions.)

\section{Step 2: Passing to the quotients}

Recall $\coh\PP^n\simeq \grmod S/\grmod_{\mathrm{fl}}S$, the quotient by finite-length (torsion) modules (Serre), so
\[
\Dcal^b(\coh\PP^n)\simeq \Dcal^b(\grmod S)/\Dcal^b_{\mathrm{fl}}(\grmod S).
\]
Finite-length modules are built from copies of $\k(d)$, and $\mathbf{R}(\k)$ is free over $E$. Therefore $\mathbf{R}$ carries $\Dcal^b_{\mathrm{fl}}(\grmod S)$ onto the thick subcategory $\perf(E)$ of perfect complexes, and Koszul duality descends to
\[
\Dcal^b(\coh\PP^n)\simeq \Dcal^b(\grmod E)/\perf(E).
\]
Finally, since $E$ is Frobenius, $\grmod E$ is a Frobenius category, and the natural functor $\stgrmod E\to \Dcal^b(\grmod E)/\perf(E)$ is an equivalence (Rickard, Buchweitz). This proves the theorem.

\section{Step 3: Explicit form via Tate resolutions}

For $\mathcal F\in\coh\PP^n$ the equivalence can be made explicit. Let $M=\bigoplus_d H^0(\mathcal F(d))$ be its module of twisted sections. Applying $\mathbf{R}$ to high truncations $M_{\ge r}$ gives an exact complex of free $E$-modules in high degrees; extending to the left by a free resolution yields the \emph{Tate resolution} $\mathbf T(\mathcal F)$, an acyclic complex of free $E$-modules with terms
\[
\mathbf T(\mathcal F)^e=\bigoplus_i \omega_E(e-i)\otimes H^i(\mathcal F(e-i)).
\]
Thus $\mathbf T(\mathcal F)$ packages all cohomology groups of all twists of $\mathcal F$. Syzygy modules of $\mathbf T(\mathcal F)$ are the objects of $\stgrmod E$ corresponding to $\mathcal F$ (up to shift), and $\mathbf T$ is the functor realizing the equivalence of the theorem on objects.

\section{Step 4: Beilinson's resolution of the diagonal}

On $\PP^n\times\PP^n$ there is the exact sequence
\[
0\to \mathcal O(-n)\boxtimes\Omega^n(n)\to\cdots\to\mathcal O(-1)\boxtimes\Omega^1(1)\to \mathcal O\boxtimes\mathcal O\to\mathcal O_\Delta\to0 .
\]
It yields Beilinson's monad: every $\mathcal F$ is the cohomology of a complex with terms $\bigoplus_j H^i(\mathcal F(-j))\otimes\Omega^j(j)$, whose differentials are read off from $\mathbf T(\mathcal F)$. In particular $\mathcal O(-n),\dots,\mathcal O$ is a full exceptional collection, so $\Dcal^b(\coh\PP^n)$ is generated by finitely many explicit objects. This is the geometric counterpart of the algebraic generation used in Step~1, and gives an alternative route to full faithfulness and essential surjectivity: check on the generators $\Omega^j(j)$, which correspond to the simple modules $\k$ (suitably shifted) over $E$, then extend by d\'evissage.

\section{Summary of the strategy}

\begin{enumerate}
\item \textbf{Koszul duality}: $\Dcal^b(\grmod S)\simeq\Dcal^b(\grmod E)$, proved on generators ($S(d)\leftrightarrow\k$) and extended by generation.
\item \textbf{Quotients}: torsion $S$-modules correspond to perfect $E$-complexes, so Serre localization on the $S$-side matches Verdier localization on the $E$-side.
\item \textbf{Frobenius property}: $\Dcal^b(\grmod E)/\perf(E)\simeq\stgrmod E$.
\item \textbf{Explicit inverses}: Tate resolutions and Beilinson monads encode sheaf cohomology and recover $\mathcal F$ from the $E$-module.
\end{enumerate}

\begin{remark}
Gradings, shifts, and which of $V$, $V^*$ appears vary between sources, so exact twists above should be checked against a reference. A readable account is Eisenbud--Fl{\o}ystad--Schreyer, \emph{Sheaf cohomology and free resolutions over exterior algebras}, Trans.\ AMS 355 (2003).
\end{remark}

\end{document}
